% from ega4-18-fr.tex
\oldpage[IV-4]{135}

Il est immédiat que $d)$ implique la condition $b)$ de (18.5.11) car une immersion fermée étale est une immersion ouverte (18.9.1). Inversement, supposons vérifiées les conditions de (18.5.11) et prouvons $d)$. L'hypothèse de $d)$ implique que $f$ est quasi-fini au point $x$ (Err$_{\mathrm{III}}$, 20), en vertu de (13.1.4). Donc en vertu de la condition $c)$ de (18.5.11), $\sh{O}_{X,x}$ est une $A$-algèbre finie et $\Spec(\sh{O}_{X,x})$ un voisinage ouvert $U$ de $x$ dans $X$. En outre, comme $\sh{O}_{X,x}/\mathfrak m_s\sh{O}_{X,x}$ est isomorphe à $k=A/\mathfrak m_s$, le lemme de Nakayama prouve que l'homomorphisme $A\to\sh{O}_{X,x}$ est surjectif, donc que $f|U$ est une immersion fermée.

\begin{proposition}[18.5.19]
\label{IV.18.5.19}
Soient $A$ un anneau local noethérien et hensélien, $S=\Spec(A)$, $s$ le point fermé de $S$, $f:X\to S$ un morphisme propre; posons $X_s=f^{-1}(s)$. Alors l'application $Y\longmapsto Y\cap X_s$ est une bijection de l'ensemble des composantes connexes de $X$ sur l'ensemble des composantes connexes de $X_s$.
\end{proposition}

En considérant un sous-préschéma fermé de $X$ ayant pour espace sous-jacent une composante connexe de $X$, on est ramené à prouver que si $X$ est \emph{connexe et non vide}, alors $X_s$ est \emph{connexe et non vide}. Le fait que $X_s$ soit non vide résulte de (II, 7.2.1); pour prouver que $X_s$ est connexe, raisonnons par l'absurde, en considérant la factorisation de Stein $f:X\xrightarrow{f'}S'\xrightarrow{g}S$ du morphisme propre $f$ (III, 4.3.3); par hypothèse, l'ensemble discret fini $g^{-1}(s)$ contiendrait au moins deux points. Puisque $A$ est hensélien et $g$ séparé et de type fini, il résulterait alors de (18.5.11, $c)$ que $S'$ serait somme de deux préschémas non vides $S'_1$, $S'_2$ (puisque l'intersection de l'un d'eux avec $g^{-1}(s)$ est réduite à un seul point). Puisque $f'$ est surjectif (III, 4.3.1), on en conclurait que $X$ est somme de deux préschémas non vides, contrairement à l'hypothèse.

\subsection{Hensélisation.}
\EGAReaderTarget{ega4.c04.subsec.18-6}{subsection:IV.18.6}

\begin{env}[18.6.1]
\EGAReaderTarget{ega4.c04.statement.18-6-1}{IV.18.6.1}
Étant donné un anneau local $A$, nous dirons qu'une $A$-algèbre locale $B$ est \emph{essentiellement étale} s'il existe une $A$-algèbre étale $C$ et un idéal premier $\mathfrak n$ de $C$ tels que $B$ soit $A$-isomorphe à $C_{\mathfrak n}$ et que l'homomorphisme composé $A\to C\to C_{\mathfrak n}$ soit \emph{local} (autrement dit, que $\mathfrak n$ soit au-dessus de l'idéal maximal $\mathfrak m$ de $A$); cela entraîne (17.6.1) que $\mathfrak n$ est le \emph{seul} idéal premier de $B$ au-dessus de $\mathfrak m$ puisque $B$ est un anneau local; si $B$, $B'$ sont deux $A$-algèbres locales essentiellement étales, tout $A$-homomorphisme $B\to B'$ est donc \emph{local}.

Si $A'$ est une $A$-algèbre locale essentiellement étale et $A''$ une $A'$-algèbre locale essentiellement étale, alors $A''$ est une $A$-algèbre locale essentiellement étale. En effet, on a par hypothèse $A'=B_{\mathfrak n}$, où $B$ est une $A$-algèbre étale et $\mathfrak n$ un idéal premier de $B$ au-dessus de l'idéal maximal de $A$, et $A''=B'_{\mathfrak n'}$, où $B'$ est une $A'$-algèbre étale et $\mathfrak n'$ un idéal premier de $B'$ au-dessus de l'idéal $\mathfrak nB_{\mathfrak n}$. Si l'on pose $S=B-\mathfrak n$, de sorte que $A'=S^{-1}B$, $B'$ est de la forme $S^{-1}C$, où $C$ est une $B$-algèbre de présentation finie. Par suite $C$ est une $A$-algèbre de présentation finie et $A''$ est de la forme $C_{\mathfrak r}$, où $\mathfrak r$ est un idéal premier de $C$ au-dessus de l'idéal maximal de $A$. Comme $A'$ est une $A$-algèbre formellement étale et $A''$ une $A'$-algèbre formellement étale, $A''$ est une $A$-algèbre formellement

